\section{可微流形}\label{sec:differentiable-manifolds}

回顾一下, $n$ 维拓扑流形是具有可数开集基的 Hausdorff 空间, 其中每个点都有一个与 $\mathbb{R}^n$ 的某个开子集同胚的邻域. 设 $U$ 是 $\mathbb{R}^n$ 的一个开子集, $\varphi:U\to\varphi(U)\subset M$ 是从 $U$ 到 $M$ 的一个开子集的同胚, 则每个这样的二元组 $(U,\varphi)$ 都称为一个\textbf{参数化}. 其逆映射 $\varphi^{-1}$ 称为\textbf{坐标系}或\textbf{坐标图}, 集合 $\varphi(U)\subset M$ 称为\textbf{坐标邻域}. 当两个坐标邻域相交时, 我们可以写出相应的坐标变换 (参见 Fig.~\ref{fig:chart-transition-maps}). 得到可微流形的想法是选取一部分参数化, 使得坐标变换都是可微映射.

\RGfigure[.57\textwidth]{chart-transition-maps}{参数化与重叠映射}

\begin{RiemannDefinition}\label{def:differentiable-manifold}
$n$ 维\textbf{可微流形}或\textbf{光滑流形}由一个 $n$ 维拓扑流形 $M$ 和一族定义在开集 $U_\alpha\subset\mathbb{R}^n$ 上的参数化 $\varphi_\alpha:U_\alpha\to M$ 组成, 且满足:
\begin{enumerate}[label=(\roman*),leftmargin=*]
\item 坐标邻域覆盖 $M$, 即 $\bigcup_\alpha\varphi_\alpha(U_\alpha)=M$;
\item 对每一对满足
\[
  W:=\varphi_\alpha(U_\alpha)\cap\varphi_\beta(U_\beta)\ne\varnothing,
\]
的指标 $\alpha,\beta$, 重叠映射
\begin{gather*}
  \varphi_\beta^{-1}\circ\varphi_\alpha:\varphi_\alpha^{-1}(W)\to\varphi_\beta^{-1}(W)\\
  \varphi_\alpha^{-1}\circ\varphi_\beta:\varphi_\beta^{-1}(W)\to\varphi_\alpha^{-1}(W)
\end{gather*}
都是 $C^\infty$ 的;
\item 族 $\mathcal{A}=\{(U_\alpha,\varphi_\alpha)\}$ 关于 (i) 和 (ii) 是极大的, 即如果 $\varphi_0:U_0\to M$ 是一个参数化, 使得对 $\mathcal{A}$ 中的所有 $\varphi$, $\varphi_0^{-1}\circ\varphi$ 和 $\varphi^{-1}\circ\varphi_0$ 都是 $C^\infty$ 的, 那么 $(U_0,\varphi_0)$ 属于 $\mathcal{A}$.
\end{enumerate}
\end{RiemannDefinition}

\begin{RiemannRemark}\label{rem:maximal-and-equivalent-atlases}
\begin{enumerate}[label=(\arabic*),leftmargin=*]
\item 任何满足 (i) 和 (ii) 的族 $\mathcal{A}=\{(U_\alpha,\varphi_\alpha)\}$ 都称为 $M$ 的一个 \textbf{$C^\infty$ 图册}. 如果 $\mathcal{A}$ 还满足 (iii), 就称其为\textbf{极大图册}或\textbf{可微结构}.
\item 条件 (iii) 纯粹是技术性的. 给定 $M$ 上任意图册 $\mathcal{A}=\{(U_\alpha,\varphi_\alpha)\}$, 存在唯一的包含它的极大图册 $\widetilde{\mathcal{A}}$. 实际上, 可以取 $\widetilde{\mathcal{A}}$ 为与 $\mathcal{A}$ 中每个参数化都满足 (ii) 的所有参数化所成的集合. 显然 $\mathcal{A}\subset\widetilde{\mathcal{A}}$, 且容易检验 $\widetilde{\mathcal{A}}$ 满足 (i) 和 (ii). 由构造可知, $\widetilde{\mathcal{A}}$ 关于 (i) 和 (ii) 是极大的. 如果两个图册定义了\phantomsection\label{idx:equivalent-atlases}同一个可微结构, 就称它们\textbf{等价}.
\item 我们也可以通过要求坐标变换是 $C^k$ 映射来定义 $C^k$ 流形 (此时 $C^0$ 流形就表示拓扑流形).
\end{enumerate}
\end{RiemannRemark}

\begin{RiemannExample}\label{ex:examples-of-differentiable-structures}
\begin{enumerate}[label=(\arabic*),ref=\arabic*,leftmargin=*]
\item\label{ex:standard-atlas-on-euclidean-space} 赋有 Euclidean 度量所定义的通常拓扑的空间 $\mathbb{R}^n$ 是 Hausdorff 空间, 并且具有可数的开集基. 例如, 若只考虑一个参数化 $(\mathbb{R}^n,id)$, 则 Definition~\ref{def:differentiable-manifold} 的 (i) 和 (ii) 自然满足, 从而得到 $\mathbb{R}^n$ 的一个图册. 包含这个参数化的极大图册通常称为 $\mathbb{R}^n$ 上的\textbf{标准可微结构}. 当然也可以考虑其他图册. 例如, 取由参数化 $(\mathbb{R}^n,\varphi)$ 定义的图册, 其中 $\varphi(x)=Ax$, $A$ 是非奇异的 $(n\times n)$ 矩阵. 证明这两个图册等价是一道简单的练习.
\item\label{ex:inequivalent-atlases-on-the-line} 一个流形可能有不等价的图册. 考虑 $\mathbb{R}$ 上的两个图册 $\{(\mathbb{R},\varphi_1)\}$ 和 $\{(\mathbb{R},\varphi_2)\}$, 其中 $\varphi_1(x)=x$, $\varphi_2(x)=x^3$. 由于映射 $\varphi_2^{-1}\circ\varphi_1$ 在原点不可微, 这两个图册定义了不同的可微结构 (不过, 我们将看到, 它们彼此微分同胚) [参见 Exercise~\ref{ex:exercise-differentiable-structures}(\ref{ex:prove-inequivalence-of-line-atlases}) 和 Exercise~\ref{ex:properties-of-differentiable-maps}(\ref{ex:properties-of-differentiable-maps-part-six})].
\item\label{ex:open-subsets-inherit-atlases} 光滑流形的每个开子集 $V$ 都是同维数的流形. 实际上, 由于 $V$ 是 $M$ 的子集, 它的子空间拓扑是 Hausdorff 的, 且具有可数的开集基. 此外, 若 $\mathcal{A}=\{(U_\alpha,\varphi_\alpha)\}$ 是 $M$ 的图册, 取所有满足 $\varphi_\alpha(U_\alpha)\cap V\ne\varnothing$ 的 $U_\alpha$, 容易检验参数化族 $\widetilde{\mathcal{A}}=\{(\widetilde U_\alpha,\varphi_\alpha|_{\widetilde U_\alpha})\}$ 是 $V$ 的一个图册, 其中 $\widetilde U_\alpha=\varphi_\alpha^{-1}(V)$.
\item\label{ex:matrix-space-and-general-linear-group} 设 $M_{n\times n}$ 是实系数 $n\times n$ 矩阵的集合. 把矩阵元素排成一行, 可见这个空间就是 $\mathbb{R}^{n^2}$, 因而它是流形. 由第~\ref{ex:open-subsets-inherit-atlases}~个例子可知, $GL(n)=\{A\in M_{n\times n}\mid\det A\ne 0\}$ 也是一个 $n^2$ 维流形. 实际上, 行列式是从 $M_{n\times n}$ 到 $\mathbb{R}$ 的连续映射, 而 $GL(n)$ 是开集 $\mathbb{R}\setminus\{0\}$ 的原像.
\item\label{ex:sphere-atlas-by-coordinate-projections} 考虑 \textbf{$n$ 维球面}
\[
  S^n=\{(x^1,\ldots,x^{n+1})\in\mathbb{R}^{n+1}\mid(x^1)^2+\cdots+(x^{n+1})^2=1\}
\]
以及映射
\[
  \begin{aligned}
  \varphi_i^+&:U\subset\mathbb{R}^n\to S^n\\
  (x^1,\ldots,x^n)&\mapsto(x^1,\ldots,x^{i-1},g(x^1,\ldots,x^n),x^i,\ldots,x^n),
  \end{aligned}
\]
\[
  \begin{aligned}
  \varphi_i^-&:U\subset\mathbb{R}^n\to S^n\\
  (x^1,\ldots,x^n)&\mapsto(x^1,\ldots,x^{i-1},-g(x^1,\ldots,x^n),x^i,\ldots,x^n),
  \end{aligned}
\]
其中
\[
  U=\{(x^1,\ldots,x^n)\in\mathbb{R}^n\mid (x^1)^2+\cdots+(x^n)^2<1\}
\]
且
\[
  g(x^1,\ldots,x^n)=\bigl(1-(x^1)^2-\cdots-(x^n)^2\bigr)^{\frac12}.
\]
球面是 $\mathbb{R}^{n+1}$ 的子集, 因而 (赋予子空间拓扑后) 是 Hausdorff 空间, 且具有可数的开集基. 也容易检验族 $\{(U,\varphi_i^+),(U,\varphi_i^-)\}_{i=1}^{n+1}$ 是 $S^n$ 的图册, 因而这个空间是 $n$ 维流形 (相应的坐标图恰是到超平面 $x^i=0$ 的投影).
\item 可以为立方体的表面 $Q\subset\mathbb{R}^3$ 定义一个图册, 使它成为光滑流形. 假定立方体以原点为中心, 考虑由 $f(x)=x/\lVert x\rVert$ 定义的映射 $f:Q\to S^2$. 于是, 对 $S^2$ 的图册 $\{(U_\alpha,\varphi_\alpha)\}$, 族 $\{(U_\alpha,f^{-1}\circ\varphi_\alpha)\}$ 定义了 $Q$ 的一个图册.
\end{enumerate}
\end{RiemannExample}

\begin{RiemannRemark}\label{rem:topological-manifolds-without-smooth-structure}
有些拓扑流形完全不容许任何可微结构. 实际上, Kervaire 在 1960 年给出了第一个例子 (一个 10 维流形) \RGcite{Ker60}, 不久 Smale 又构造了另一个例子 (维数为 12) \RGcite{Sma60}. 早在 1956 年, Milnor \RGcite{Mil07} 就给出了一个他认为不容许可微结构的 8 维流形的例子, 但这一点直到 1965 年才得到证明 (见 \RGcite{Nov65}).
\end{RiemannRemark}

\begin{RiemannExercise}\label{ex:exercise-differentiable-structures}
\begin{enumerate}[label=(\arabic*),ref=\arabic*,leftmargin=*]
\item\label{ex:compare-atlases-by-union} 证明光滑流形的两个图册 $\mathcal{A}_1$ 和 $\mathcal{A}_2$ 等价, 当且仅当 $\mathcal{A}_1\cup\mathcal{A}_2$ 是图册.
\item\label{ex:open-sets-in-parameterizations} 设 $M$ 是可微流形. 证明集合 $V\subset M$ 是开集, 当且仅当对于一个 $C^\infty$ 图册中的每个参数化 $(U_\alpha,\varphi_\alpha)$, $\varphi_\alpha^{-1}(V)$ 都是 $\mathbb{R}^n$ 的开子集.
\item\label{ex:euclidean-atlases-equivalent} 证明 Example~\ref{ex:examples-of-differentiable-structures}(\ref{ex:standard-atlas-on-euclidean-space}) 中 $\mathbb{R}^n$ 上的两个图册等价.
\item\label{ex:prove-inequivalence-of-line-atlases} 考虑 Example~\ref{ex:examples-of-differentiable-structures}(\ref{ex:inequivalent-atlases-on-the-line}) 中 $\mathbb{R}$ 上的两个图册 $\{(\mathbb{R},\varphi_1)\}$ 和 $\{(\mathbb{R},\varphi_2)\}$, 其中 $\varphi_1(x)=x$, $\varphi_2(x)=x^3$. 证明 $\varphi_2^{-1}\circ\varphi_1$ 在原点不可微, 并由此得出这两个图册不等价.
\item\label{ex:surfaces-as-local-graphs} 回顾初等向量微积分中的定义: 曲面 $S\subset\mathbb{R}^3$ 是一个集合, 对每个 $p\in S$, 在 $\mathbb{R}^3$ 中存在 $p$ 的一个邻域 $V_p$, 并存在一个 $C^\infty$ 映射 $f_p:U_p\to\mathbb{R}$ (其中 $U_p$ 是 $\mathbb{R}^2$ 的开子集), 使得 $S\cap V_p$ 是 $z=f_p(x,y)$, $x=f_p(y,z)$ 或 $y=f_p(x,z)$ 的图像. 证明 $S$ 是 2 维光滑流形.
\item\label{ex:product-differentiable-manifold} (乘积流形) 设 $\{(U_\alpha,\varphi_\alpha)\}$ 和 $\{(V_\beta,\psi_\beta)\}$ 分别为两个光滑流形 $M$ 和 $N$ 的图册. 证明族 $\{(U_\alpha\times V_\beta,\varphi_\alpha\times\psi_\beta)\}$ 是乘积 $M\times N$ 的一个图册. 赋予由这个图册生成的可微结构后, $M\times N$ 称为 $M$ 和 $N$ 的\textbf{乘积流形}.
\item\label{ex:stereographic-sphere-atlas} (球极投影) 考虑带子空间拓扑的 $n$ 维球面 $S^n$, 设 $N=(0,\ldots,0,1)$ 和 $S=(0,\ldots,0,-1)$ 分别为北极和南极. 从 $N$ 出发的\textbf{球极投影}是映射 $\pi_N:S^n\setminus\{N\}\to\mathbb{R}^n$, 它把点 $p\in S^n\setminus\{N\}$ 送到经过 $N$ 和 $p$ 的直线与超平面 $x^{n+1}=0$ 的交点 (参见 Fig.~\ref{fig:stereographic-projection}). 类似地, 从 $S$ 出发的球极投影是映射 $\pi_S:S^n\setminus\{S\}\to\mathbb{R}^n$, 它把 $S^n\setminus\{S\}$ 上的点 $p$ 送到经过 $S$ 和 $p$ 的直线与同一超平面的交点. 验证 $\{(\mathbb{R}^n,\pi_N^{-1}),(\mathbb{R}^n,\pi_S^{-1})\}$ 是 $S^n$ 的一个图册. 证明这个图册与 Example~\ref{ex:examples-of-differentiable-structures}(\ref{ex:sphere-atlas-by-coordinate-projections}) 中的图册等价. 由这两个图册得到的极大图册称为 $S^n$ 上的\textbf{标准可微结构}.
\item\label{ex:projective-space-atlas-from-sphere} (\emph{实射影空间}) \textbf{实射影空间} $\mathbb{R}P^n$ 是 $\mathbb{R}^{n+1}$ 中过原点的直线所成的集合. 这个空间可定义为 $S^n$ 关于等价关系 $x\sim -x$ 的商空间, 即将每一点与它的对径点等同.
\begin{enumerate}[label=(\alph*),leftmargin=*]
\item 证明赋予商拓扑的商空间 $\mathbb{R}P^n=S^n/{\sim}$ 是 Hausdorff 空间, 且具有可数的开集基. (\textbf{Hint:} 使用 Proposition~\RGref{prop:topology-of-quotients}.)
\item 考虑 Example~\ref{ex:examples-of-differentiable-structures}(\ref{ex:sphere-atlas-by-coordinate-projections}) 中定义的 $S^n$ 上的图册, 以及由 $\pi(x)=[x]$ 给出的典范投影 $\pi:S^n\to\mathbb{R}P^n$, 为 $\mathbb{R}P^n$ 定义一个图册.
\end{enumerate}
\item\label{ex:projective-space-atlas-from-lines} 也可以用另一种方式在 $\mathbb{R}P^n$ 上定义图册, 即将它等同于 $\mathbb{R}^{n+1}\setminus\{0\}$ 关于等价关系 $x\sim\lambda x$ (其中 $\lambda\in\mathbb{R}\setminus\{0\}$) 的商空间. 为此, 考虑集合 $V_i=\{[x^1,\ldots,x^{n+1}]\mid x^i\ne 0\}$ (对应于 $\mathbb{R}^{n+1}$ 中过原点而不包含在超平面 $x^i=0$ 内的直线所成的集合), 并考虑由
\[
  \varphi_i(x^1,\ldots,x^n)=[x^1,\ldots,x^{i-1},1,x^i,\ldots,x^n].
\]
定义的映射 $\varphi_i:\mathbb{R}^n\to V_i$. 证明:
\begin{enumerate}[label=(\alph*),leftmargin=*]
\item 族 $\{(\mathbb{R}^n,\varphi_i)\}$ 是 $\mathbb{R}P^n$ 的一个图册;
\item 这个图册与 Exercise~\ref{ex:exercise-differentiable-structures}(\ref{ex:projective-space-atlas-from-sphere}) 中的图册定义相同的可微结构.
\end{enumerate}
\item\label{ex:non-hausdorff-line-with-two-origins} (\emph{非 Hausdorff 流形}) 设 $M$ 是 $\mathbb{R}$ 与一点 $p$ 的不交并. 考虑映射 $f_i:\mathbb{R}\to M$ ($i=1,2$), 定义为当 $x\in\mathbb{R}\setminus\{0\}$ 时 $f_i(x)=x$, 并且 $f_1(0)=0$, $f_2(0)=p$. 证明:
\begin{enumerate}[label=(\alph*),leftmargin=*]
\item 映射 $f_i^{-1}\circ f_j$ 在其定义域上可微;
\item 若考虑由 $\{(\mathbb{R},f_1),(\mathbb{R},f_2)\}$ 构成的图册, 则相应的拓扑不满足 Hausdorff 公理.
\end{enumerate}
\end{enumerate}
\end{RiemannExercise}

\RGfigure[.57\textwidth]{stereographic-projection}{球极投影}
